\begin{table*}[t]
\centering
\small
\caption{\textbf{Example of a knowledge-matched contrast.}
An A2 and an A3 contrast from the same seed. The options and the gold answer (A) are the same for both: A) adsorption, B) headspace method, C) solvent extraction, D) steam distillation.}
\label{tab:example_contrast}
\begin{tabular}{@{}p{0.07\textwidth}p{0.43\textwidth}p{0.43\textwidth}@{}}
\toprule
 & \textbf{Korean} & \textbf{English translation} \\
\midrule
Anchor (excerpt)
& 숯(Charcoal)이나 활성탄을 시료 용기의 상부공간에 두어 휘발성 촉진제 성분을 표면에 붙잡아 포집하는 방식은 흡착법이다.
\newline 상부공간법은 시료 위의 기체층인 헤드스페이스 자체를 채취하거나 분석하여 휘발성 성분을 확인하는 방법이다.
& Placing charcoal or activated carbon in the headspace of the sample container, so that volatile accelerant components are captured on its surface, is the adsorption method.
\newline The headspace method identifies volatile components by collecting or analyzing the headspace itself, the gas layer above the sample. \\
\midrule
A2
& 화재 잔류물이 담긴 밀폐병을 실험실로 옮긴 뒤, 분석자는 병마개 안쪽에 작은 검은 탄소성 알갱이가 든 금속망을 매달았다. 병을 따뜻하게 유지하여 휘발 성분이 시료 위 빈 공간으로 퍼지게 한 후, 금속망만 꺼내 분석했더니 휘발성 촉진제 성분이 확인되었다. 이 촉진제 회수 방법으로 가장 적절한 것은?
& After a sealed jar containing fire debris was moved to the laboratory, the analyst suspended a metal mesh holding small black carbonaceous granules inside the lid. The jar was kept warm so that volatile components spread into the empty space above the sample, and when only the mesh was taken out and analyzed, volatile accelerant components were identified. Which method of recovering the accelerant is most appropriate? \\
\midrule
A3
& 화재 잔류물이 든 금속 용기를 온수 욕조에 올려 둔 채 밀봉하였다. 뚜껑 안쪽에는 작은 활성탄 조각을 실로 매달아 잔류물 위 빈 공간에 놓이게 했고, 시간이 지나자 뚜껑 안쪽에는 물방울도 조금 맺혔다. 분석자는 꺼낸 활성탄 조각을 소량의 용매가 든 바이알에 넣은 뒤 그 액을 장비에 주입하였다. 이 절차에서 촉진제 성분을 회수한 방법은 무엇인가?
& A metal container holding fire debris was sealed while placed in a hot water bath. A small piece of activated carbon was hung by a thread inside the lid so that it sat in the empty space above the debris, and after some time a few water droplets also formed inside the lid. The analyst placed the removed piece of activated carbon in a vial containing a small amount of solvent and injected the liquid into the instrument. By which method was the accelerant recovered in this procedure? \\
\bottomrule
\end{tabular}
\end{table*}